%!TEX root = ../sztuthesis_main.tex
% 设置中文摘要


\addcontentsline{toc}{section}{摘要}

\keywordscn{智能代码审查；LLM；检索增强生成（RAG）；Agent}
\categorycn{TP391}
\begin{abstractcn}
### 摘要  
随着软件开发规模与复杂度的持续增长，代码质量和安全性已成为影响项目成功的关键因素。传统代码审查依赖人工经验，存在效率低下、覆盖不全等问题，难以应对现代软件的快速迭代需求。为此，本文设计并实现了一个基于大语言模型（LLM）与检索增强生成（RAG）技术的智能代码审查VS Code插件系统，旨在提升代码审查的自动化水平与准确性。

系统采用前后端分离架构：前端为VS Code插件，负责用户交互与代码分析触发；后端为Agent服务，集成静态分析工具、Tree-sitter解析器及CodeLlama大语言模型，实现多维度代码审查与智能修复。核心功能包括：基于规则与AI的混合审查机制、实时审查与批量审查、智能代码修复建议、基于RAG的上下文增强分析等。

在实现过程中，系统整合了多种代码分析技术（如Pylint、ESLint）与向量数据库（ChromaDB），通过RAG技术增强LLM对代码上下文的理解，同时引入反思机制优化审查结果。测试结果表明，系统能够有效识别代码中的安全漏洞、性能问题与风格缺陷，生成的修复代码与原代码功能一致且质量显著提升。

本研究验证了大语言模型与RAG技术在代码审查场景的应用价值，为开发人员提供了一种高效、智能的代码质量保障工具。未来可通过扩展支持更多编程语言、优化模型推理速度，进一步提升系统的实用性与普适性。

% 图X幅，表X个，参考文献X篇（四号宋体）

\end{abstractcn}